\chapter{Lymph Node Biopsy}

\section*{Overview}
A lymph node biopsy removes all or part of an enlarged or abnormal lymph node for histologic and molecular diagnosis, most commonly to evaluate suspected lymphoma or metastatic cancer. The exact approach, setting, and preparation depend on the indication, anatomy, and the patient's health.

\section*{Alternative Names}
Excisional lymph node biopsy; sentinel lymph node biopsy; lymphadenectomy; node sampling

\section*{Principle}
The lymph node is excised intact (excisional biopsy) or sampled (incisional or core biopsy) and examined by histology, flow cytometry, and molecular testing. Excisional biopsy provides the architecture needed for lymphoma classification.

\section*{History}
Lymph node biopsy became a cornerstone of pathology in the 19th century, and excisional biopsy remains the standard for suspected lymphoma, complemented by fine-needle and core biopsy in modern practice.

\section*{Why It Is Done}
Performed for unexplained persistent lymphadenopathy, suspected lymphoma or metastatic disease, and for sentinel node evaluation in breast cancer, melanoma, and other cancers.

\section*{Is This a Primary or Secondary Test or Procedure}
Primary diagnostic procedure for lymphadenopathy of unknown cause.

\section*{How It Is Performed}
Performed under local or general anesthesia. The node is localized by palpation or ultrasound, an incision is made along skin creases, and the node is carefully dissected and removed intact with hemostasis. The wound is closed and the specimen is oriented and sent fresh for processing.

\section*{Preparation}
Usually none beyond routine preoperative assessment. The surgeon and pathologist coordinate handling to ensure fresh tissue for flow cytometry and molecular studies.

\section*{Contraindications and Precautions}
Bleeding disorders and infection over the site are relative contraindications. Precaution: identify and protect adjacent structures such as vessels and nerves, especially in the neck and axilla.

\section*{Risks}
Bleeding, infection, seroma, nerve injury (e.g., spinal accessory nerve in the neck), lymphedema, and anesthetic risks.

\section*{Aftercare and Recovery}
Most patients go home the same day. Mild pain is controlled with analgesics, and the wound is kept clean and dry.

\section*{Outcome and Follow-up}
Histology distinguishes reactive hyperplasia from lymphoma, metastasis, granulomatous disease, or infection; flow cytometry and cytogenetics subclassify lymphomas.

\section*{Expected Results}
Reactive lymphoid hyperplasia with a preserved architecture and no evidence of malignancy.

\section*{Follow-up}
Results guide staging, referral, and treatment; persistent or recurring adenopathy elsewhere prompts further evaluation.

\section*{When to Seek Medical Care}
Follow the procedure team's specific instructions. Seek urgent medical assessment for severe or worsening pain, uncontrolled bleeding, fever or signs of infection, trouble breathing, chest pain, fainting, new weakness or confusion, inability to urinate, or any rapidly worsening symptom. The appropriate warning signs depend on the procedure and the patient's medical condition.

\section*{References}
\begin{enumerate}
  \item MedlinePlus. Lymph node biopsy. U.S. National Library of Medicine; 2025.
  \item Current Medical Diagnosis and Treatment. McGraw-Hill; 2025.
\end{enumerate}

\clearpage
